\newcommand{\nomeEstado}{BAHIA}
\newcommand{\siglaEstado}{BA}
\newcommand{\nomeConcessionaria}{NEOENERGIA BAHIA}
\newcommand{\cnpjConcessionaria}{15.139.629/0001-94}
\newcommand{\termoCIP}{Sexto}
\newcommand{\presidenteConcessionaria}{Sr. Thiago Guth.}

\newcommand{\empresaResponsavel}{ABEL}
\newcommand{\nomeMunicipio}{Jaguaripe}
\newcommand{\nomeMunicipioMaisculo}{\MakeUppercase{\nomeMunicipio}}
\newcommand{\IPestimada}{2748133}
\newcommand{\cnpjMunicipio}{13.796.289/0001-49}
\newcommand{\telefone}{(75) 3642-2112}
\newcommand{\email}{contato@jaguaripe.ba.gov.br}
\newcommand{\sedePrefeitura}{Praça Histórica, nº 01, Centro, Jaguaripe – BA, CEP 44480-000}
\newcommand{\nomePrefeito}{Fábio Nonato Barbosa}
\newcommand{\generoPrefeitura}{pelo Prefeito}
\newcommand{\numeroContrato}{166/2025}
\newcommand{\quantidadeAditivos}{1}
\newcommand{\legitimidade}{1}
\newcommand{\publicacao}{Diário Oficial do Município}
\newcommand{\dataPublicacao}{24 de Abril de 2025}
\newcommand{\tituloClausula}{Cláusula Primeira}
\newcommand{\clausula}{[...]

Deverá a empresa efetuar  levantamento dos créditos a que faz jus o município, referentes aos pagamentos indevidos à  concessionária de energia elétrica, conforme os prazos estabelecidos na Resolução Normativa  da ANEEL nº 1.000, de 7 de dezembro de 2021, especialmente o art. 323, 82º, inciso I, e suas  alterações posteriores, podendo para tanto, representar administrativamente o Município de  Jaguaripe - BA perante a Agência Nacional de Energia Elétrica — ANEEL, em todos os  processos, reclamações, recursos e demais atos decorrentes de suas competências.}
